% !TeX program = lualatex
% !TeX root = projekt 1.tex
% ============================================================
% Zajednički preamble za sva predavanja RWA kolegija
% Uključiti s: % ============================================================
% Zajednički preamble za sva predavanja RWA kolegija
% Uključiti s: % ============================================================
% Zajednički preamble za sva predavanja RWA kolegija
% Uključiti s: \input{../common/preamble.tex}
% ============================================================

\documentclass[aspectratio=149,11pt]{beamer}
\usetheme{CambridgeUS}

\usepackage{fontspec}              % radi samo sa XeLaTeX/LuaLaTeX
\setsansfont{Latin Modern Sans}    % sans-serif font
\setmonofont{Latin Modern Mono}    % monospace font za kod
\usepackage[croatian]{babel}
\usepackage{csquotes}              % \enquote{} → hrvatski navodnici „..."
\usepackage{lmodern}
\usepackage{booktabs}
\usepackage{listings}
\usepackage{xcolor}
\usepackage{hyperref}

% ---------- Autor i logo ----------
\author{mag.ing. Josip Torić (SIT UNIZD)}
\titlegraphic{\includegraphics[height=2.5cm]{../common/SIT_logo_tekst1.pdf}}

% ---------- Listings ----------
\definecolor{codebg}{RGB}{248,248,248}
\definecolor{codeframeclr}{RGB}{220,220,220}
\lstdefinelanguage{bash}{
  morekeywords={sudo,apt,brew,pip,pipx,uv,winget,python,git,flask,source,deactivate,code,gunicorn,docker,compose,heroku,render,railway,node,npm,pnpm,yarn,uvicorn,alembic,pytest,ruff,black,precommit},
  sensitive=false,
  morecomment=[l]{\#}
}
\lstdefinelanguage{python}{
  morekeywords={from,import,def,class,return,if,elif,else,for,in,while,try,except,raise,with,as,not,and,or,is,lambda,yield,True,False,None,await,async},
  sensitive=true,
  morecomment=[l]{\#},
  morestring=[b]',
  morestring=[b]"
}
\lstdefinelanguage{sql}{
  morekeywords={SELECT,FROM,WHERE,INSERT,INTO,VALUES,UPDATE,SET,DELETE,CREATE,TABLE,PRIMARY,KEY,FOREIGN,REFERENCES,UNIQUE,NOT,NULL,INDEX,ON,ALTER,DROP,CONSTRAINT},
  sensitive=false,
  morecomment=[l]{--}
}
\lstdefinelanguage{json}{
  morecomment=[l]{//},
  morestring=[b]",
  showstringspaces=false
}
\lstdefinelanguage{typescript}{
  morekeywords={import,export,from,const,let,var,function,return,if,else,for,of,in,while,switch,case,break,async,await,try,catch,finally,throw,new,class,this,extends,implements,interface,type,enum,readonly,public,private,protected,static,get,set,true,false,null,undefined,void,string,number,boolean,any,unknown,never,typeof,keyof,as},
  sensitive=true,
  morecomment=[l]{//},
  morecomment=[s]{/*}{*/},
  morestring=[b]',
  morestring=[b]",
  morestring=[b]`
}
\lstalias{javascript}{typescript}
\lstset{
  backgroundcolor=\color{codebg},
  basicstyle=\ttfamily\footnotesize,
  frame=single,
  rulecolor=\color{codeframeclr},
  frameround=tttt,
  breaklines=true,
  tabsize=2,
  showstringspaces=false,
  upquote=true,
  columns=fullflexible
}

% ---------- Helper macros ----------
\newcommand{\codefile}[1]{\texttt{#1}}
\newcommand{\kbd}[1]{\texttt{#1}}

% ---------- Code frame environment ----------
% Automatski dodaje fragile i shrink=5, koristi se umjesto \begin{frame}
% Primjer: \begin{codeframe}{Naslov} ... \end{codeframe}
% Primjer s opcijama: \begin{codeframe}[shrink=15]{Naslov} ... \end{codeframe}
\newenvironment{codeframe}[2][]
  {\begin{frame}[fragile,shrink=5,#1]{#2}}
  {\end{frame}}

% ============================================================

\documentclass[aspectratio=149,11pt]{beamer}
\usetheme{CambridgeUS}

\usepackage{fontspec}              % radi samo sa XeLaTeX/LuaLaTeX
\setsansfont{Latin Modern Sans}    % sans-serif font
\setmonofont{Latin Modern Mono}    % monospace font za kod
\usepackage[croatian]{babel}
\usepackage{csquotes}              % \enquote{} → hrvatski navodnici „..."
\usepackage{lmodern}
\usepackage{booktabs}
\usepackage{listings}
\usepackage{xcolor}
\usepackage{hyperref}

% ---------- Autor i logo ----------
\author{mag.ing. Josip Torić (SIT UNIZD)}
\titlegraphic{\includegraphics[height=2.5cm]{../common/SIT_logo_tekst1.pdf}}

% ---------- Listings ----------
\definecolor{codebg}{RGB}{248,248,248}
\definecolor{codeframeclr}{RGB}{220,220,220}
\lstdefinelanguage{bash}{
  morekeywords={sudo,apt,brew,pip,pipx,uv,winget,python,git,flask,source,deactivate,code,gunicorn,docker,compose,heroku,render,railway,node,npm,pnpm,yarn,uvicorn,alembic,pytest,ruff,black,precommit},
  sensitive=false,
  morecomment=[l]{\#}
}
\lstdefinelanguage{python}{
  morekeywords={from,import,def,class,return,if,elif,else,for,in,while,try,except,raise,with,as,not,and,or,is,lambda,yield,True,False,None,await,async},
  sensitive=true,
  morecomment=[l]{\#},
  morestring=[b]',
  morestring=[b]"
}
\lstdefinelanguage{sql}{
  morekeywords={SELECT,FROM,WHERE,INSERT,INTO,VALUES,UPDATE,SET,DELETE,CREATE,TABLE,PRIMARY,KEY,FOREIGN,REFERENCES,UNIQUE,NOT,NULL,INDEX,ON,ALTER,DROP,CONSTRAINT},
  sensitive=false,
  morecomment=[l]{--}
}
\lstdefinelanguage{json}{
  morecomment=[l]{//},
  morestring=[b]",
  showstringspaces=false
}
\lstdefinelanguage{typescript}{
  morekeywords={import,export,from,const,let,var,function,return,if,else,for,of,in,while,switch,case,break,async,await,try,catch,finally,throw,new,class,this,extends,implements,interface,type,enum,readonly,public,private,protected,static,get,set,true,false,null,undefined,void,string,number,boolean,any,unknown,never,typeof,keyof,as},
  sensitive=true,
  morecomment=[l]{//},
  morecomment=[s]{/*}{*/},
  morestring=[b]',
  morestring=[b]",
  morestring=[b]`
}
\lstalias{javascript}{typescript}
\lstset{
  backgroundcolor=\color{codebg},
  basicstyle=\ttfamily\footnotesize,
  frame=single,
  rulecolor=\color{codeframeclr},
  frameround=tttt,
  breaklines=true,
  tabsize=2,
  showstringspaces=false,
  upquote=true,
  columns=fullflexible
}

% ---------- Helper macros ----------
\newcommand{\codefile}[1]{\texttt{#1}}
\newcommand{\kbd}[1]{\texttt{#1}}

% ---------- Code frame environment ----------
% Automatski dodaje fragile i shrink=5, koristi se umjesto \begin{frame}
% Primjer: \begin{codeframe}{Naslov} ... \end{codeframe}
% Primjer s opcijama: \begin{codeframe}[shrink=15]{Naslov} ... \end{codeframe}
\newenvironment{codeframe}[2][]
  {\begin{frame}[fragile,shrink=5,#1]{#2}}
  {\end{frame}}

% ============================================================

\documentclass[aspectratio=149,11pt]{beamer}
\usetheme{CambridgeUS}

\usepackage{fontspec}              % radi samo sa XeLaTeX/LuaLaTeX
\setsansfont{Latin Modern Sans}    % sans-serif font
\setmonofont{Latin Modern Mono}    % monospace font za kod
\usepackage[croatian]{babel}
\usepackage{csquotes}              % \enquote{} → hrvatski navodnici „..."
\usepackage{lmodern}
\usepackage{booktabs}
\usepackage{listings}
\usepackage{xcolor}
\usepackage{hyperref}

% ---------- Autor i logo ----------
\author{mag.ing. Josip Torić (SIT UNIZD)}
\titlegraphic{\includegraphics[height=2.5cm]{../common/SIT_logo_tekst1.pdf}}

% ---------- Listings ----------
\definecolor{codebg}{RGB}{248,248,248}
\definecolor{codeframeclr}{RGB}{220,220,220}
\lstdefinelanguage{bash}{
  morekeywords={sudo,apt,brew,pip,pipx,uv,winget,python,git,flask,source,deactivate,code,gunicorn,docker,compose,heroku,render,railway,node,npm,pnpm,yarn,uvicorn,alembic,pytest,ruff,black,precommit},
  sensitive=false,
  morecomment=[l]{\#}
}
\lstdefinelanguage{python}{
  morekeywords={from,import,def,class,return,if,elif,else,for,in,while,try,except,raise,with,as,not,and,or,is,lambda,yield,True,False,None,await,async},
  sensitive=true,
  morecomment=[l]{\#},
  morestring=[b]',
  morestring=[b]"
}
\lstdefinelanguage{sql}{
  morekeywords={SELECT,FROM,WHERE,INSERT,INTO,VALUES,UPDATE,SET,DELETE,CREATE,TABLE,PRIMARY,KEY,FOREIGN,REFERENCES,UNIQUE,NOT,NULL,INDEX,ON,ALTER,DROP,CONSTRAINT},
  sensitive=false,
  morecomment=[l]{--}
}
\lstdefinelanguage{json}{
  morecomment=[l]{//},
  morestring=[b]",
  showstringspaces=false
}
\lstdefinelanguage{typescript}{
  morekeywords={import,export,from,const,let,var,function,return,if,else,for,of,in,while,switch,case,break,async,await,try,catch,finally,throw,new,class,this,extends,implements,interface,type,enum,readonly,public,private,protected,static,get,set,true,false,null,undefined,void,string,number,boolean,any,unknown,never,typeof,keyof,as},
  sensitive=true,
  morecomment=[l]{//},
  morecomment=[s]{/*}{*/},
  morestring=[b]',
  morestring=[b]",
  morestring=[b]`
}
\lstalias{javascript}{typescript}
\lstset{
  backgroundcolor=\color{codebg},
  basicstyle=\ttfamily\footnotesize,
  frame=single,
  rulecolor=\color{codeframeclr},
  frameround=tttt,
  breaklines=true,
  tabsize=2,
  showstringspaces=false,
  upquote=true,
  columns=fullflexible
}

% ---------- Helper macros ----------
\newcommand{\codefile}[1]{\texttt{#1}}
\newcommand{\kbd}[1]{\texttt{#1}}

% ---------- Code frame environment ----------
% Automatski dodaje fragile i shrink=5, koristi se umjesto \begin{frame}
% Primjer: \begin{codeframe}{Naslov} ... \end{codeframe}
% Primjer s opcijama: \begin{codeframe}[shrink=15]{Naslov} ... \end{codeframe}
\newenvironment{codeframe}[2][]
  {\begin{frame}[fragile,shrink=5,#1]{#2}}
  {\end{frame}}


\title[Razvoj web aplikacija]{Razvoj web aplikacija \\ Projekt 1 - Backend}
\date{05.10.2026.}

\begin{document}

\frame{\titlepage}

% ==============================
% DIO 1: PROJEKT I BODOVANJE
% ==============================

% ------------------------------
\begin{frame}{Plan prezentacije}
\begin{enumerate}
  \item Što gradite i kako se boduje
  \item Detaljni zahtjevi za backend
  \item Bodovna raspodjela 1. dijela
  \item Instalacija razvojnog okruženja
  \item GitHub Education, Copilot i Claude Code
  \item Kako koristiti AI da stvarno naučite
  \item Što vas čeka na 1. obrani
\end{enumerate}
\end{frame}

% ------------------------------
\begin{frame}{Projekt kroz semestar}
Tijekom semestra izrađujete \textbf{vlastitu full-stack web aplikaciju}. Tema je vaš izbor, ali mora zadovoljiti zahtjeve iz ove prezentacije.

\vspace{0.6em}
\begin{tabular}{@{}lp{5.2cm}lr@{}}
  \toprule
  \textbf{Dio} & \textbf{Što se predaje} & \textbf{Obrana} & \textbf{Bodovi} \\
  \midrule
  1. dio & Backend: API, baza, testovi & 16.11. & 45 \\
  2. dio & Frontend SPA + deploy & 18.01. & 45 \\
  \multicolumn{2}{@{}l}{Dolasci na predavanja} & & 10 \\
  \midrule
  \multicolumn{3}{@{}l}{\textbf{Ukupno}} & \textbf{100} \\
  \bottomrule
\end{tabular}

\vspace{0.6em}
Svaki dio se brani zasebno, usmeno, uz demonstraciju i pitanja o kodu. Danas govorimo o 1. dijelu - zahtjevi za frontend dolaze u prezentaciji \enquote{Projekt 2} na početku frontend bloka.
\end{frame}

% ------------------------------
\begin{frame}{Tema projekta: kako odabrati}
Dobra tema ima \textbf{uloge, pravila i statuse} - ne samo tablice koje se pune i čitaju.

\vspace{0.4em}
\begin{itemize}
  \item Najmanje dvije role s različitim ovlastima: admin i korisnik, organizator i sudionik, trener i sportaš
  \item Podaci koji pripadaju nekome - temelj za ownership provjere
  \item Barem jedno pravilo koje ovisi o vremenu, statusu ili kapacitetu
  \item Realan proces koji možete objasniti u dvije rečenice
\end{itemize}

\vspace{0.6em}
Ako imate samo CRUD bez pravila i rola - to nije dovoljno za prolaz. Ako niste sigurni je li tema dovoljna, pitajte me \textbf{prije} nego što počnete.
\end{frame}

% ------------------------------
\begin{frame}{Primjeri tema}
\begin{columns}[T]
\column{0.48\textwidth}
\begin{itemize}
  \item Sustav rezervacija termina
  \item Event management s prijavama
  \item Task / workflow sustav
  \item Coach \& athlete tracker
\end{itemize}

\column{0.48\textwidth}
\begin{itemize}
  \item Document approval sustav
  \item Mini marketplace s oglasima
  \item Evidencija članstva u klubu
  \item Sustav narudžbi i dostave
\end{itemize}
\end{columns}

\vspace{0.6em}
Referentni primjer s predavanja je sustav prijava na powerlifting natjecanje: admin i klubovi, natjecatelji, rokovi nominacija. Vaša tema mora imati usporedivu složenost.
\end{frame}

% ------------------------------
\begin{frame}{Tehnologije}
Koristimo konkretan stack - isti za sve studente. Na predavanjima učimo svaku tehnologiju od nule.

\vspace{0.4em}
\begin{tabular}{@{}lp{7cm}@{}}
  \toprule
  \textbf{Tehnologija} & \textbf{Uloga} \\
  \midrule
  Python 3.11+ & jezik backenda \\
  FastAPI & web framework \\
  SQLAlchemy 2.0 & ORM, async \\
  Alembic & migracije baze \\
  PostgreSQL & baza podataka \\
  Pydantic & validacija i DTO \\
  pytest & testovi \\
  uv & paketi i virtualna okruženja \\
  Docker Compose & baza u kontejneru \\
  \bottomrule
\end{tabular}

\vspace{0.4em}
Drugi stack nije dopušten - cilj kolegija je dubina na jednom stacku, ne širina na pet.
\end{frame}

% ------------------------------
\begin{frame}{Struktura repozitorija}
Projekt je \textbf{monorepo}: backend i frontend u istom repozitoriju na GitHubu.

\vspace{0.4em}
\begin{tabular}{@{}lp{7.5cm}@{}}
  \toprule
  \textbf{Folder} & \textbf{Sadržaj} \\
  \midrule
  \codefile{api/routers/} & HTTP endpointi, bez poslovne logike \\
  \codefile{api/services/} & poslovna logika i provjere \\
  \codefile{api/repositories/} & svi upiti prema bazi \\
  \codefile{api/models/} & SQLAlchemy ORM definicije \\
  \codefile{api/schemas/} & Pydantic DTO modeli \\
  \codefile{api/core/} & config, session, auth, errori \\
  \codefile{api/tests/} & pytest testovi \\
  \codefile{web/} & frontend - dolazi u 2. dijelu \\
  \bottomrule
\end{tabular}
\end{frame}

% ------------------------------
\begin{frame}{Arhitektura: slojevi su obavezni}
\vspace{0.2em}
\begin{center}
\texttt{Routers $\rightarrow$ Services $\rightarrow$ Repositories $\rightarrow$ Models $\rightarrow$ PostgreSQL}
\end{center}

\vspace{0.6em}
Svaki sloj ima jednu odgovornost:
\begin{itemize}
  \item Router prima request, parsira parametre, poziva service, vraća response
  \item Service sadrži poslovnu logiku - ne zna za HTTP ni SQL
  \item Repository radi sve upite prema bazi i vraća model objekte
\end{itemize}

\vspace{0.6em}
Ako vam je sva logika u routeru ili router direktno piše SQL - to se mora popraviti \textbf{prije} obrane. Arhitektura nije preporuka, nego uvjet.
\end{frame}

% ==============================
% DIO 2: DETALJNI ZAHTJEVI
% ==============================

% ------------------------------
\begin{frame}{Zahtjev 1: baza podataka i migracije}
\begin{itemize}
  \item \textbf{Najmanje 3 tablice} s relacijama: 1:N obavezno, M:N ako tema traži
  \item Svaka tablica ima surrogate \codefile{id} kao primarni ključ
  \item \textbf{Constrainti na razini baze}: NOT NULL za obavezna polja, UNIQUE gdje domena traži jedinstvenost, FOREIGN KEY za sve veze
  \item Sve promjene sheme kroz \textbf{Alembic migracije} - ručni SQL na bazi nije dopušten
  \item Svaka migracija ima funkcionalan \codefile{downgrade()}
  \item \textbf{Seed skripta}: idempotentna, kreira role i demo podatke za obranu
\end{itemize}

\vspace{0.4em}
Na obrani provjeravam: \enquote{Zašto je ovaj stupac UNIQUE? Što se dogodi ako constraint makneš i dva requesta stignu istovremeno?}
\end{frame}

% ------------------------------
\begin{frame}{Zahtjev 2: autentikacija}
\begin{itemize}
  \item Login s \codefile{username} + \codefile{password}
  \item Lozinke hashirane \textbf{bcryptom} - plaintext ili MD5/SHA su nula bodova za autentikaciju
  \item \textbf{JWT access token} (15 min) + \textbf{refresh token} (dulji rok)
  \item Tri endpointa: \codefile{/auth/login}, \codefile{/auth/refresh}, \codefile{/auth/me}
  \item Kriva lozinka i nepostojeći username vraćaju \textbf{isti} 401 error - bez otkrivanja koji korisnici postoje
  \item \codefile{JWT\_SECRET} iz env varijable, nikad hardkodiran u kodu
\end{itemize}

\vspace{0.4em}
JWT payload sadrži samo ono što autorizacija treba: \codefile{sub}, \codefile{role}, vlasnički ID, \codefile{exp}. Nikad lozinku ni osobne podatke.
\end{frame}

% ------------------------------
\begin{frame}{Zahtjev 3: autorizacija i ownership}
\begin{itemize}
  \item \textbf{Najmanje dvije role} s različitim ovlastima, npr. admin i korisnik
  \item Zaštita rola kroz dependency (\codefile{require\_role}), ne if-ovima po routerima
  \item \textbf{Ownership}: korisnik vidi i mijenja \textbf{samo svoje} resurse - provjera u service layeru
  \item Ispravni status kodovi: 401 bez tokena, 403 s tokenom ali bez prava
\end{itemize}

\vspace{0.6em}
\begin{tabular}{@{}lll@{}}
  \toprule
  \textbf{Scenarij} & \textbf{Status} & \textbf{Zašto} \\
  \midrule
  Nema tokena ili istekao & 401 & Ne znamo tko si \\
  Kriva rola & 403 & Znamo tko si, ne smiješ \\
  Tuđi resurs & 403 ili 404 & Ovisno o politici otkrivanja \\
  \bottomrule
\end{tabular}

\vspace{0.4em}
Na obrani mijenjam ID u URL-u i očekujem da vaš backend to zabrani - to je IDOR test.
\end{frame}

% ------------------------------
\begin{frame}{Zahtjev 4: poslovna pravila}
Barem \textbf{jedno pravilo koje nije CRUD} - nešto što service provjerava prije nego što dopusti akciju:

\vspace{0.4em}
\begin{tabular}{@{}lp{7.2cm}@{}}
  \toprule
  \textbf{Tip pravila} & \textbf{Primjer} \\
  \midrule
  Rok & Prijava moguća samo do datuma X \\
  Status / tranzicija & Zahtjev ide draft $\rightarrow$ poslan $\rightarrow$ odobren, ne preskače korake \\
  Kapacitet & Termin prima najviše N rezervacija \\
  Jedinstvenost u kontekstu & Isti korisnik se ne može dvaput prijaviti na isti događaj \\
  \bottomrule
\end{tabular}

\vspace{0.6em}
Pravilo živi u \textbf{service layeru} i vraća domenski error s jasnim kodom, npr. \codefile{deadline\_passed} uz status 400 ili 409. UI smije sakriti gumb, ali backend \textbf{mora} zabraniti.
\end{frame}

% ------------------------------
\begin{frame}{Zahtjev 5: API surface}
Najmanje \textbf{10-15 endpointova}, organiziranih oko resursa:

\vspace{0.4em}
\begin{tabular}{@{}llr@{}}
  \toprule
  \textbf{Grupa} & \textbf{Primjer} & \textbf{Min.} \\
  \midrule
  Auth & login, refresh, me & 3 \\
  Glavni entitet & CRUD: list, create, get, update, delete & 4-5 \\
  Drugi entitet & CRUD s ownership provjerom & 4-5 \\
  Workflow & akcije: odobri, odbij, odjavi, zaključaj & 1-2 \\
  \bottomrule
\end{tabular}

\vspace{0.6em}
\begin{itemize}
  \item \textbf{Nested URL-ovi} gdje ima smisla: \codefile{/clubs/\{id\}/lifters}, ne \codefile{/lifters?club\_id=5}
  \item Lista endpointova ima \textbf{paginaciju} (\codefile{limit}/\codefile{offset})
  \item Ispravni status kodovi: 201 create, 204 delete, 400 pravilo, 403 zabranjeno, 404 ne postoji, 409 duplikat, 422 validacija
\end{itemize}
\end{frame}

% ------------------------------
\begin{frame}{Zahtjev 6: validacija i error handling}
\begin{itemize}
  \item \textbf{Pydantic scheme za svaki input}: zasebne Create, Update i Response varijante - nikad isti DTO za input i output
  \item Response scheme ne izlažu interne podatke: \codefile{hashed\_password} nikad ne izlazi iz API-ja
  \item Validatori za domenska ograničenja: datum u prošlosti, dopuštene vrijednosti, rokovi u ispravnom poretku
  \item \textbf{Standardiziran error format}: svaki error vraća JSON s \codefile{code} i \codefile{message}, kroz globalni exception handler
  \item Partial update kroz PATCH s \codefile{exclude\_unset=True}
\end{itemize}

\vspace{0.4em}
422 vraća Pydantic za oblik podataka, 400/409 vraća service za smisao. Na obrani pitam u čemu je razlika.
\end{frame}

% ------------------------------
\begin{frame}{Zahtjev 7: testovi}
Pytest testovi koji dokazuju da pravila vrijede:

\vspace{0.4em}
\begin{tabular}{@{}lp{7cm}@{}}
  \toprule
  \textbf{Područje} & \textbf{Minimum} \\
  \midrule
  Autentikacija & uspješan login, kriva lozinka, istekli token, refresh \\
  Autorizacija & svaka rola: što smije i što dobiva 403 \\
  Ownership & korisnik A ne vidi i ne mijenja podatke korisnika B \\
  Poslovna pravila & pravilo prolazi, pravilo zabranjuje, granični slučaj \\
  Constrainti & duplikat vraća 409, nevažeći input 422 \\
  \bottomrule
\end{tabular}

\vspace{0.6em}
Svaki endpoint ima barem jedan \textbf{negativni test}. Vremenski ovisna pravila testirajte \textbf{freezegunom} - test ne smije ovisiti o danu pokretanja.
\end{frame}

% ------------------------------
\begin{frame}{Zahtjev 8: Git i predaja}
\begin{itemize}
  \item Projekt je na \textbf{GitHubu} od prvog tjedna - repozitorij mi šaljete odmah
  \item \textbf{Commitovi kroz semestar}, ne jedan veliki commit dan prije obrane - povijest commitova je dio ocjene kvalitete
  \item Smislene commit poruke: \enquote{feat: add registration deadline check}, ne \enquote{update}
  \item \codefile{.env} nikad u gitu - u repozitoriju stoji \codefile{.env.example}
  \item \codefile{README.md} s uputama za pokretanje: baza, migracije, seed, server
\end{itemize}

\vspace{0.4em}
Repozitorij s tri commita od kojih je jedan \enquote{final version} govori mi više od bilo kojeg odgovora na usmenom.
\end{frame}

% ------------------------------
\begin{frame}{Bodovna raspodjela 1. dijela (45 bodova)}
\begin{tabular}{@{}lr@{}}
  \toprule
  \textbf{Komponenta} & \textbf{Bodovi} \\
  \midrule
  Modeliranje baze, migracije i seed & 6 \\
  Autentikacija: bcrypt, JWT, refresh & 7 \\
  Autorizacija: role i ownership & 7 \\
  Poslovna pravila & 7 \\
  API dizajn i status kodovi & 5 \\
  Validacija i error handling & 4 \\
  Testovi & 5 \\
  Arhitektura i kvaliteta koda & 4 \\
  \midrule
  \textbf{Ukupno} & \textbf{45} \\
  \bottomrule
\end{tabular}

\vspace{0.4em}
Bodovi vrijede samo uz \textbf{razumijevanje}: dio koji na usmenom ne znate objasniti ne nosi bodove, ma koliko dobro radio.
\end{frame}

% ------------------------------
\begin{frame}{Checklist za predaju 1. dijela}
Prije obrane provjerite da sve prolazi na svježem cloneu repozitorija:

\vspace{0.4em}
\begin{enumerate}
  \item \codefile{docker compose up -d db} podiže bazu
  \item \codefile{alembic upgrade head} prolazi bez greške
  \item \codefile{python -m app.seed} puni demo podatke
  \item \codefile{uvicorn app.main:app} se podiže, \codefile{/docs} pokazuje sve endpointe
  \item \codefile{pytest} je zelen
  \item README objašnjava sve gornje korake
\end{enumerate}

\vspace{0.4em}
\enquote{Radi na mom laptopu} nije predaja. Ako ja ne mogu pokrenuti projekt po README-u, obrana počinje s minusom.
\end{frame}

% ==============================
% DIO 3: INSTALACIJA
% ==============================

% ------------------------------
\begin{frame}{Što trebate instalirati}
\begin{tabular}{@{}lll@{}}
  \toprule
  \textbf{Alat} & \textbf{Verzija} & \textbf{Svrha} \\
  \midrule
  Git & najnovija & verzioniranje koda \\
  Python & 3.11+ & backend (FastAPI) \\
  Docker Desktop & najnovija & PostgreSQL baza \\
  VS Code & najnovija & editor \\
  \bottomrule
\end{tabular}

\vspace{0.6em}
Node.js za frontend dolazi u 2. dijelu - za backend vam ne treba. Sve se instalira kroz službene stranice ili package managere; detalji na sljedećim slideovima.
\end{frame}

% ------------------------------
\begin{frame}[fragile,shrink=5]{Git}
\begin{lstlisting}[language=bash]
# Windows - preuzmite installer sa https://git-scm.com
# macOS
brew install git
# Linux (Debian/Ubuntu)
sudo apt install git

# Provjera
git --version

# Osnovna konfiguracija (obavezno nakon instalacije)
git config --global user.name "Ime Prezime"
git config --global user.email "email@example.com"
\end{lstlisting}

\vspace{0.4em}
Ako nemate GitHub račun, napravite ga odmah - trebat će vam za predaju projekta i za GitHub Education.
\end{frame}

% ------------------------------
\begin{frame}[fragile,shrink=5]{Python i uv}
\begin{lstlisting}[language=bash]
# Windows - preuzmite sa https://python.org
# OBAVEZNO oznacite "Add Python to PATH" u installeru

# macOS
brew install python@3.12

# Linux
sudo apt install python3

# uv - moderni alat za pakete i virtualna okruzenja
# Windows: winget install astral-sh.uv
# macOS:   brew install uv
# Linux:   curl -LsSf https://astral.sh/uv/install.sh | sh

# Provjera
python --version   # ili python3 --version
uv --version
\end{lstlisting}

\vspace{0.4em}
Virtualna okruženja i pakete vodimo kroz \textbf{uv}: \texttt{uv venv} stvara okruženje, \texttt{uv pip install} instalira pakete - brže od klasičnog pipa, a svaki projekt ostaje izoliran od sustava.
\end{frame}

% ------------------------------
\begin{frame}[fragile,shrink=5]{Docker Desktop}
Preuzmite Docker Desktop s \texttt{docker.com}. Uključuje Docker Engine i Docker Compose.

\begin{lstlisting}[language=bash]
# Provjera nakon instalacije
docker --version
docker compose version

# Pokretanje PostgreSQL baze za razvoj
docker compose up -d db
\end{lstlisting}

\vspace{0.4em}
Docker koristimo \textbf{samo za bazu} - ne trebate kontejnerizirati svoj backend. PostgreSQL lokalno bez Dockera je moguć, ali Docker Compose je jednostavniji jer konfiguracija ostaje u \codefile{docker-compose.yml}.
\end{frame}

% ------------------------------
\begin{frame}{VS Code i ekstenzije}
\textbf{VS Code} preuzmite s \texttt{code.visualstudio.com}. Preporučene ekstenzije za backend:

\vspace{0.4em}
\begin{itemize}
  \item \textbf{Python} (Microsoft) - IntelliSense, debugging, venv podrška
  \item \textbf{SQLTools} + PostgreSQL driver - pregled baze iz editora
  \item \textbf{REST Client} ili \textbf{Thunder Client} - testiranje API-ja bez Postmana
\end{itemize}

\vspace{0.6em}
Smijete koristiti i drugi editor (PyCharm, IntelliJ), ali upute na predavanjima prate VS Code.
\end{frame}

% ==============================
% DIO 4: AI ALATI
% ==============================

% ------------------------------
\begin{frame}{GitHub Education - što dobivate}
Kao studenti imate pravo na \textbf{GitHub Student Developer Pack} koji uključuje:

\begin{itemize}
  \item \textbf{GitHub Copilot} besplatno - AI autocomplete u editoru
  \item Pristup \textbf{AI modelima} kroz Copilot Chat - najnovije generacije GPT i Claude modela
  \item GitHub Pro - privatni repozitoriji, Actions minute, Codespaces
  \item Razni bonusi: domene, cloud krediti, alati
\end{itemize}

\vspace{0.4em}
Aktivacija: \texttt{education.github.com} $\rightarrow$ prijavite se s GitHub računom $\rightarrow$ potvrdite studentski status (index ili potvrda s fakulteta).
\end{frame}

% ------------------------------
\begin{frame}{GitHub Copilot u VS Code-u}
Nakon aktivacije GitHub Education, instalirajte \textbf{GitHub Copilot} ekstenziju u VS Code.

\vspace{0.4em}
\textbf{Dva načina korištenja:}
\begin{enumerate}
  \item \textbf{Inline suggestions} - dok tipkate, Copilot predlaže nastavak koda. Tab za prihvaćanje, Esc za odbijanje.
  \item \textbf{Copilot Chat} (Ctrl+Shift+I) - razgovor s AI-jem o kodu. Selektirajte dio koda i pitajte za objašnjenje, refactoring ili debugging.
\end{enumerate}

\vspace{0.4em}
U Copilot Chatu birate model - isprobajte različite, svaki ima drugačije snage.
\end{frame}

% ------------------------------
\begin{frame}{Copilot Chat - korisne naredbe}
\begin{tabular}{@{}ll@{}}
  \toprule
  \textbf{Naredba} & \textbf{Što radi} \\
  \midrule
  \texttt{/explain} & Objašnjava selektirani kod \\
  \texttt{/fix} & Predlaže popravak za grešku \\
  \texttt{/tests} & Generira testove za selektirani kod \\
  \texttt{/doc} & Generira dokumentaciju \\
  \texttt{@workspace} & Daje kontekst cijelog projekta \\
  \bottomrule
\end{tabular}

\vspace{0.6em}
Selektirajte funkciju $\rightarrow$ otvorite Chat $\rightarrow$ napišite \texttt{/explain} ili postavite pitanje prirodnim jezikom. Chat vidi vaš kod i daje kontekstualni odgovor.
\end{frame}

% ------------------------------
\begin{frame}{Claude Code u VS Code-u}
\textbf{Claude Code} je Anthropicov agentski alat koji se integrira u VS Code kao ekstenzija.

\vspace{0.4em}
\textbf{Instalacija:}
\begin{enumerate}
  \item Napravite račun na \texttt{claude.ai}
  \item U VS Code instalirajte ekstenziju \textbf{Claude Code} (Anthropic)
  \item Prijavite se kroz ekstenziju - otvori se browser za autentikaciju
  \item Claude se pojavljuje kao panel u VS Code-u
\end{enumerate}

\vspace{0.6em}
Claude čita vaše datoteke, razumije strukturu projekta i može predlagati promjene direktno u editoru. Za razliku od Copilota koji radi inline, Claude radi u \textbf{agentskom} modu - može čitati, pisati i izvršavati naredbe.
\end{frame}

% ------------------------------
\begin{frame}{Claude vs Copilot - kad koristiti koji}
\begin{columns}[T]
\column{0.48\textwidth}
\textbf{Copilot}
\begin{itemize}
  \item Brzi inline prijedlozi dok tipkate
  \item Autocomplete za poznate obrasce
  \item Kratka pitanja o kodu (\texttt{/explain})
  \item Generiranje boilerplate koda
\end{itemize}

\column{0.48\textwidth}
\textbf{Claude Code}
\begin{itemize}
  \item Složeniji zadaci s više datoteka
  \item Debugging - čita error, traži uzrok
  \item Razumijevanje arhitekture
  \item Pisanje i pokretanje testova
\end{itemize}
\end{columns}

\vspace{0.6em}
Ne trebate birati samo jedan. Copilot za svakodnevno tipkanje, Claude kad trebate razmisliti o problemu ili napraviti veću promjenu.
\end{frame}

% ------------------------------
\begin{frame}{AI smijete koristiti - uz razumijevanje}
AI alati su dio modernog razvoja softvera i zabraniti ih nema smisla. Cilj kolegija nije provjeriti možete li pisati kod iz glave, nego \textbf{razumijete li što taj kod radi i zašto}.

\vspace{0.6em}
Na usmenom ne gledam \textbf{jeste li} koristili AI. Gledam \textbf{razumijete li} kod koji ste predali. Ako ne možete objasniti vlastiti kod, svejedno je tko ga je napisao - bodove ne nosi.
\end{frame}

% ------------------------------
\begin{frame}{Kako AI koristiti za učenje}
\textbf{Pristup koji radi:}
\begin{enumerate}
  \item \textbf{Prvo pokušajte sami} - razmislite o problemu 10-15 minuta prije nego što pitate AI
  \item \textbf{Pitajte za objašnjenje, ne za gotov kod} - \enquote{Objasni mi kako radi JWT refresh flow} umjesto \enquote{Napiši mi refresh endpoint}
  \item \textbf{Čitajte generirani kod red po red} - što ne razumijete, tražite da objasni
  \item \textbf{Reproducirajte sami} - nakon što razumijete, napišite sličan kod bez AI-ja
  \item \textbf{Testirajte rubne slučajeve} - što ako token istekne? ako korisnik nema pravo?
\end{enumerate}

\vspace{0.4em}
\textbf{Što ne raditi}: copy-paste cijelog modula bez čitanja, traženje kompletnog featurea odjednom, slijepo ponavljanje \enquote{ne radi, popravi} umjesto čitanja error poruke.
\end{frame}

% ------------------------------
\begin{frame}{Dobar prompt vs loš prompt}
\begin{columns}[T]
\column{0.48\textwidth}
\textbf{Loš prompt}
\begin{itemize}
  \item \enquote{Napravi mi login}
  \item \enquote{Ne radi, popravi}
  \item \enquote{Napiši testove}
  \item \enquote{Dodaj autorizaciju}
\end{itemize}

\column{0.48\textwidth}
\textbf{Dobar prompt}
\begin{itemize}
  \item \enquote{Objasni zašto login vraća 401 na ispravan password - evo requesta i responsea}
  \item \enquote{Imam ovaj service - objasni kako ownership check radi}
  \item \enquote{Napiši test koji provjerava da korisnik ne može vidjeti tuđe podatke}
\end{itemize}
\end{columns}

\vspace{0.6em}
Što precizniji prompt, to bolji odgovor. Dajte kontekst: error poruku, dio koda, što ste pokušali.
\end{frame}

% ==============================
% DIO 5: OBRANA
% ==============================

% ------------------------------
\begin{frame}{1. obrana (16.11.): kako izgleda}
Demonstracija API-ja kroz Swagger, pregled repozitorija i ključnih dijelova koda, pitanja o odlukama i edge caseovima.

\vspace{0.4em}
\textbf{Tipična pitanja:}
\begin{itemize}
  \item Objasni arhitekturu - što radi router, što service, što repository?
  \item Kako radi JWT flow? Gdje se hashira lozinka?
  \item Gdje provjeravaš ownership? Pokaži mi liniju koda.
  \item Kako backend zabranjuje akciju nakon isteka roka?
  \item Zašto ovaj endpoint vraća 403, a ne 404?
  \item Što testira ovaj test i što bi se pokvarilo da ga nema?
\end{itemize}

\vspace{0.4em}
Na kraju: \textbf{mini promjena uživo} - dodaj polje na model, napravi migraciju, prilagodi schemu i endpoint.
\end{frame}

% ------------------------------
\begin{frame}{Što je dovoljan odgovor na usmenom}
Ne tražim savršene odgovore. Tražim da \textbf{razumijete} svoj kod.

\vspace{0.6em}
\begin{columns}[T]
\column{0.48\textwidth}
\textbf{Prolazi}
\begin{itemize}
  \item \enquote{Dependency čita JWT iz headera, dekodira ga i vraća usera - endpoint ga dobiva kroz Depends}
  \item \enquote{403 jer je autenticiran, ali nema pravo na taj resurs}
  \item \enquote{Test dokazuje da korisnik ne može mijenjati tuđi zapis}
\end{itemize}

\column{0.48\textwidth}
\textbf{Ne prolazi}
\begin{itemize}
  \item \enquote{To mi je AI napisao, ne znam točno}
  \item \enquote{Mislim da to radi nešto s tokenima}
  \item \enquote{Nisam siguran zašto je 403, a ne 401}
\end{itemize}
\end{columns}
\end{frame}

% ------------------------------
\begin{frame}{Rokovi i sljedeći koraci}
\begin{tabular}{@{}ll@{}}
  \toprule
  \textbf{Što} & \textbf{Kada} \\
  \midrule
  Odabir teme i repozitorij na GitHubu & do 12.10. \\
  1. obrana: backend & 16.11. \\
  Prezentacija \enquote{Projekt 2} (frontend zahtjevi) & 23.11. \\
  2. obrana: full-stack + deploy & 18.01. \\
  Rezervni termin predaje & veljački rokovi \\
  \bottomrule
\end{tabular}

\vspace{0.6em}
Počnite \textbf{ovaj tjedan}: instalirajte okruženje, odaberite temu, napravite repozitorij. Svako predavanje reproducirajte na svom projektu - tko prati ritam predavanja, backend mu je gotov prije roka.

\vspace{0.6em}
\begin{center}
\Large Pitanja?
\end{center}
\end{frame}

\end{document}
